% ======================================================================
%  CHƯƠNG 2 — KIẾN THỨC CHUẨN BỊ VÀ CÁC CÔNG TRÌNH LIÊN QUAN
%  Yêu cầu (mục 2.3): khái niệm nền tảng; khảo sát >= 5 công trình theo hướng
%  tiếp cận; phân tích điểm mạnh/yếu; xác định khoảng trống nghiên cứu.
%  Độ dài tham khảo: >= 10 trang.
% ======================================================================
\chapter{Kiến thức chuẩn bị và các công trình liên quan}
\label{ch:cosoly}

\section{Bài toán học máy tự động}
\label{sec:automl}
Một quy trình học máy có giám sát trên dữ liệu dạng bảng gồm nhiều bước liên
tiếp: tiền xử lý, trích chọn/biến đổi đặc trưng, lựa chọn thuật toán và tinh
chỉnh siêu tham số. Học máy tự động (AutoML) nhằm tự động hoá việc \emph{tìm kiếm}
trên không gian các cấu hình quy trình để đạt hiệu năng tốt nhất trong một ngân
sách cho trước. Bài toán cốt lõi của AutoML thường được hình thức hoá dưới dạng
bài toán CASH (\textit{Combined Algorithm Selection and Hyperparameter
optimization})~\cite{feurer2015autosklearn}.

\begin{dinhnghia}[Bài toán CASH]
\label{def:cash}
Cho tập thuật toán $\mathcal{A} = \{A^{(1)}, \dots, A^{(R)}\}$, mỗi thuật toán
$A^{(j)}$ có không gian siêu tham số $\Lambda^{(j)}$. Cho tập huấn luyện được chia
thành $K$ fold $\{(\mathcal{T}_i, \mathcal{V}_i)\}_{i=1}^{K}$ và một hàm mất mát
$\mathcal{L}$. Bài toán CASH là tìm cặp (thuật toán, siêu tham số)
\[
    (A^{\ast}, \lambda^{\ast}) \in
    \arg\min_{A^{(j)}\in\mathcal{A},\; \lambda\in\Lambda^{(j)}}\;
    \frac{1}{K}\sum_{i=1}^{K}
    \mathcal{L}\!\left(A^{(j)}_{\lambda}(\mathcal{T}_i),\; \mathcal{V}_i\right),
\]
trong đó $A^{(j)}_{\lambda}(\mathcal{T}_i)$ là mô hình huấn luyện trên fold
$\mathcal{T}_i$ với siêu tham số $\lambda$, và biểu thức được đánh giá trên tập
kiểm chứng $\mathcal{V}_i$.
\end{dinhnghia}

\begin{nhanxet}
Không gian tìm kiếm của Định nghĩa~\ref{def:cash} vừa rời rạc (chọn thuật toán)
vừa liên tục và có điều kiện (siêu tham số phụ thuộc thuật toán), nên rất lớn và
gồ ghề. Đây là lý do các thư viện AutoML phải dùng các chiến lược tìm kiếm hiệu
quả (tối ưu Bayes, đa độ phân giải, học tổ hợp\ldots) thay vì vét cạn.
\end{nhanxet}

\section{Kiểm chứng chéo và giao thức đánh giá công bằng}
\label{sec:cv}
Để ước lượng hiệu năng một cách không thiên lệch, dữ liệu được chia thành $K$
fold theo phương pháp kiểm chứng chéo (\textit{cross-validation}). Một phép so
sánh giữa các thư viện được gọi là \emph{công bằng} khi mọi thư viện được đánh
giá trên \emph{cùng} cách chia fold, \emph{cùng} ngân sách thời gian, \emph{cùng}
độ đo và \emph{cùng} tài nguyên tính toán~\cite{gijsbers2024amlb}. Ngoài ra, các
lần chạy thất bại (do lỗi cài đặt, hết thời gian hoặc tràn bộ nhớ) cần được ghi
nhận và phân loại, vì việc âm thầm loại bỏ chúng sẽ làm sai lệch kết quả theo
hướng có lợi cho những thư viện kém ổn định.

\section{Các độ đo đánh giá}
\label{sec:dodo-lythuyet}
Đồ án sử dụng các độ đo tiêu chuẩn theo từng loại tác vụ.

\paragraph{AUC (phân lớp nhị phân).} Diện tích dưới đường cong ROC, bằng xác suất
mô hình xếp một mẫu dương ngẫu nhiên cao hơn một mẫu âm ngẫu nhiên; giá trị càng
gần $1$ càng tốt.

\paragraph{Log-loss (phân lớp đa lớp).} Với $N$ mẫu và $C$ lớp, gọi $y_{ic}\in
\{0,1\}$ là nhãn đúng và $\hat{p}_{ic}$ là xác suất dự đoán:
\[
    \text{log-loss} = -\frac{1}{N}\sum_{i=1}^{N}\sum_{c=1}^{C}
    y_{ic}\,\log \hat{p}_{ic}.
\]

\paragraph{RMSE (hồi quy).} Với giá trị thực $y_i$ và dự đoán $\hat{y}_i$:
\[
    \text{RMSE} = \sqrt{\frac{1}{N}\sum_{i=1}^{N}\left(y_i - \hat{y}_i\right)^2}.
\]
Với AUC, giá trị lớn hơn là tốt hơn; với log-loss và RMSE, giá trị nhỏ hơn là tốt
hơn. Để thống nhất chiều ``lớn hơn là tốt hơn'' khi xếp hạng, ta có thể lấy dấu âm
của log-loss và RMSE.

\begin{vidu}[Minh hoạ cách tính RMSE]
\label{vd:rmse}
Giả sử trên một tập kiểm thử gồm $N=3$ mẫu, giá trị thực là $y = (3.0,\ 5.0,\
2.5)$ và dự đoán là $\hat{y} = (2.5,\ 5.0,\ 4.0)$. Khi đó:
\[
    \text{RMSE} = \sqrt{\tfrac{1}{3}\big[(3.0-2.5)^2 + (5.0-5.0)^2 + (2.5-4.0)^2\big]}
    = \sqrt{\tfrac{1}{3}\,(0.25 + 0 + 2.25)} = \sqrt{0.8333} \approx 0.913.
\]
Sai số chủ yếu đến từ mẫu thứ ba (lệch $1.5$ đơn vị); điều này cho thấy RMSE
nhạy với các sai lệch lớn do bình phương khuếch đại chúng.
\end{vidu}

\section{Các thư viện AutoML tiêu biểu}
\label{sec:frameworks}
\paragraph{FLAML.} Thư viện nhẹ, tập trung tối ưu tiết kiệm chi phí
(\textit{cost-frugal optimization}) và tìm kiếm đa độ phân giải, cho phép đạt kết
quả tốt trong ngân sách thời gian nhỏ~\cite{wang2021flaml}.

\paragraph{H2O AutoML.} Tự động huấn luyện nhiều họ mô hình (GBM, GLM, mạng nơ-ron,
rừng ngẫu nhiên\ldots) và kết hợp chúng bằng \textit{stacked ensemble}, hướng tới
sự tiện dụng và độ chính xác cao~\cite{ledell2020h2o}.

\paragraph{AutoGluon.} Đề cao chiến lược xếp chồng nhiều tầng (\textit{multi-layer
stacking}) và kết hợp mô hình thay vì chỉ tinh chỉnh siêu tham số, thường đạt độ
chính xác cao trên dữ liệu bảng~\cite{erickson2020autogluon}.

\section{Khảo sát các công trình liên quan}
\label{sec:conghinhlienquan}
Nhóm khảo sát các công trình liên quan và phân loại theo hướng tiếp cận. Bảng
\ref{tab:relatedwork} tóm tắt so sánh; phần sau phân tích chi tiết điểm mạnh và
hạn chế của từng hướng.

\begin{table}[H]
  \centering
  \small
  \caption{Tóm tắt các công trình liên quan theo hướng tiếp cận.}
  \label{tab:relatedwork}
  \renewcommand{\arraystretch}{1.25}
  \begin{tabularx}{\linewidth}{@{}l l X@{}}
    \toprule
    \textbf{Công trình} & \textbf{Hướng tiếp cận} & \textbf{Ý tưởng chính} \\
    \midrule
    auto-sklearn~\cite{feurer2015autosklearn} & Tối ưu Bayes + meta-learning & Giải CASH bằng tối ưu Bayes, khởi tạo bằng meta-learning và xây dựng ensemble. \\
    TPOT~\cite{olson2016tpot} & Lập trình di truyền & Tiến hoá đường ống (pipeline) học máy bằng giải thuật di truyền. \\
    FLAML~\cite{wang2021flaml} & Tối ưu tiết kiệm chi phí & Tìm kiếm đa độ phân giải, hiệu quả trong ngân sách nhỏ. \\
    H2O AutoML~\cite{ledell2020h2o} & Học tổ hợp (stacking) & Huấn luyện nhiều họ mô hình và kết hợp bằng stacked ensemble. \\
    AutoGluon~\cite{erickson2020autogluon} & Xếp chồng nhiều tầng & Ưu tiên kết hợp mô hình nhiều tầng hơn là tinh chỉnh siêu tham số. \\
    AMLB~\cite{gijsbers2024amlb} & Chuẩn đánh giá & Giao thức so sánh công bằng, tái lập, ghi nhận thất bại. \\
    \bottomrule
  \end{tabularx}
\end{table}

\subsection{Hướng tối ưu siêu tham số}
Các phương pháp như auto-sklearn~\cite{feurer2015autosklearn} và
FLAML~\cite{wang2021flaml} tập trung tìm kiếm hiệu quả trên không gian CASH.
\textit{Điểm mạnh:} nền tảng lý thuyết vững, tiết kiệm chi phí (đặc biệt là
FLAML). \textit{Hạn chế:} chất lượng phụ thuộc nhiều vào không gian tìm kiếm được
định nghĩa trước và có thể kém hơn các phương pháp tổ hợp khi ngân sách lớn.

\subsection{Hướng học tổ hợp và xếp chồng}
H2O~AutoML~\cite{ledell2020h2o} và AutoGluon~\cite{erickson2020autogluon} nhấn
mạnh việc kết hợp nhiều mô hình. \textit{Điểm mạnh:} thường đạt độ chính xác cao
nhờ tận dụng tính đa dạng của các mô hình thành phần. \textit{Hạn chế:} chi phí
huấn luyện và bộ nhớ cao hơn, mô hình kết quả khó diễn giải.

\subsection{Hướng tiến hoá}
TPOT~\cite{olson2016tpot} dùng lập trình di truyền để tiến hoá cả đường ống.
\textit{Điểm mạnh:} linh hoạt, có thể khám phá cấu trúc pipeline bất ngờ.
\textit{Hạn chế:} chi phí tính toán lớn và độ ổn định kết quả thấp hơn.

\subsection{Hướng chuẩn hoá đánh giá}
AMLB~\cite{gijsbers2024amlb} không đề xuất thuật toán mới mà cung cấp một
\emph{khung đánh giá} công bằng. \textit{Điểm mạnh:} kiểm soát chặt fold, ngân
sách, độ đo, tài nguyên và ghi nhận thất bại; là nền tảng phương pháp cho đồ án
này. \textit{Hạn chế:} thiên về dòng lệnh, cấu hình phức tạp, chưa có giao diện
trực quan và chưa hỗ trợ trực tiếp việc nạp dữ liệu tuỳ chọn từ các nguồn như
Kaggle.

\section{Khoảng trống nghiên cứu}
\label{sec:khoangtrong}
Từ khảo sát trên, nhóm nhận thấy hai khoảng trống. Thứ nhất, phần lớn nghiên cứu
tập trung \emph{cải tiến thuật toán} AutoML, trong khi khía cạnh \emph{đánh giá
công bằng và tái lập} ít được chú trọng ở cấp độ công cụ thực dụng. Thứ hai, AMLB
tuy chặt chẽ về phương pháp nhưng có rào cản sử dụng: thiếu giao diện trực quan,
chưa hỗ trợ nạp dữ liệu công khai từ Kaggle và chưa tổng hợp kết quả một cách
thân thiện. Đồ án hướng tới lấp khoảng trống này bằng cách kế thừa giao thức công
bằng của AMLB và bổ sung một hệ thống hỗ trợ nạp dữ liệu, chạy đánh giá trong môi
trường cô lập và trực quan hoá kết quả (trình bày ở Chương~\ref{ch:phuongphap}).
